% Generated draft: requires mathematical and source review.
\documentclass[11pt]{article}
\usepackage[margin=25mm]{geometry}
\usepackage{fontspec}
\setmainfont{Latin Modern Roman}
\usepackage{amsmath,amssymb,mathtools,unicode-math}
\setmathfont{Latin Modern Math}
\usepackage{enumitem,needspace,xurl,hyperref}
\setlength{\parindent}{0pt}
\setlength{\parskip}{5pt}
\setlength{\emergencystretch}{3em}
\newcommand{\N}{\mathbb{N}}
\newcommand{\Z}{\mathbb{Z}}
\newcommand{\Q}{\mathbb{Q}}
\newcommand{\R}{\mathbb{R}}
\newcommand{\C}{\mathbb{C}}
\newcommand{\F}{\mathbb{F}}
\newcommand{\A}{\mathbb{A}}
\newcommand{\PP}{\mathbb{P}}
\newcommand{\E}{\mathbb{E}}
\newcommand{\eps}{\varepsilon}
\newcommand{\dd}{\,\mathrm{d}}
\begin{document}
\section*{UnsolvedMath ID: 30002528}
\textbf{Source problem code:} catalogue entry
\subsection{Definitions and mathematical statement}
A group \(G\) is finitely presented if it admits a presentation
\[
G\cong\langle g_1,\ldots,g_n\mid r_1,\ldots,r_m\rangle,
\]
where \(n,m\) are nonnegative integers and each \(r_j\) is a word in the generators and their inverses. A Kähler manifold is a complex manifold equipped with a Hermitian metric whose associated real \((1,1)\)-form is closed. Here a Kähler group means a group isomorphic to \(\pi_1(X,x)\) for some compact connected Kähler manifold \(X\) and base point \(x\in X\).

The question is whether, for every finitely presented group \(G\), there exist a compact connected Kähler manifold \(X\), a point \(x\in X\), and an injective group homomorphism
\[
\iota:G\hookrightarrow\pi_1(X,x).
\]
The manifold may depend on \(G\); no dimension bound, projectivity assumption, or finite-index requirement on the image is imposed. The supplied statement does not explicitly define a Kähler group; the formulation above uses the standard compact-manifold convention.
\subsection{Short English statement}
Does every finitely presented group embed into the fundamental group of some compact connected Kähler manifold? The supplied record labels this question open; that status is supplied dated curation and has not been independently verified.
\subsection{Sources}
\begin{itemize}
\item Local registry: \texttt{data/unsolvedmath/problems.json}, permanent UnsolvedMath ID \(30002528\), source problem code \texttt{catalogue entry}, in the Oberwolfach Reports Open Problems collection. The supplied record assigns the year 2014 and names no individual proposer.
\item Source report: \emph{Mini-Workshop: Kähler Groups} (2014), Oberwolfach Reports, DOI 10.4171/owr/2014/09, pp.\ 531--553. Listed also as ``Oberwolfach Report 9, 2014'': \url{https://doi.org/10.4171/owr/2014/09}. The supplied background describes these reports as documenting mathematical workshops and their open-problem discussions.
\item Additional supplied literature reference: ``Liu, 2026,'' \url{https://comptes-rendus.academie-sciences.fr/mathematique/item/10.5802/crmath.811.pdf}. No fuller bibliographic information is supplied.
\item Literature assessment dated exactly \texttt{2026-08-21}: the supplied curation labels the question open and describes the embedding conjecture as an unresolved strengthening of known singular-projective realization results. This is supplied dated literature triage derived from the OWR curation, not an independently verified assessment. Time-sensitive status and cited claims require independent verification; the supplied record does not state the precise singular-projective realization results.
\item The registry labels the statement \texttt{corrected\_verified} and says its mathematical content was preserved while notation, typography, or omitted definitions were repaired for publication. This reports the registry's qualification and does not constitute independent mathematical or source verification of this draft.
\end{itemize}
\end{document}
